\chapter{Nephrotic Syndrome}
\index{Nephrotic Syndrome}
\label{sec:Nephrotic Syndrome}

\section{Overview}
Nephrotic syndrome is a clinical complex resulting from severe glomerular podocyte effacement and disruption of the charge and size barrier of the glomerular basement membrane. This genetic disorder results from specific gene mutations or chromosomal abnormalities. It may be present at birth or manifest later in life depending on the underlying mechanism. Genetic disorders result from abnormalities in an individual's DNA, ranging from single-gene mutations to chromosomal abnormalities. These conditions may be inherited from parents or arise from new mutations. Pattern of inheritance depends on whether the mutation is dominant, recessive, or X-linked. Genetic counseling helps families understand risks and make informed decisions.
\section{Causes}
This condition results from disruption of normal nephrotic syndrome function.

\begin{itemize}[noitemsep]
  \item This condition happens because of primary kidney podocytopathies or systemic disease damage to the filtration barrier. Primary etiologies include minimal change disease (most common in children), focal segmental glomerulosclerosis (FSGS, most common in African Americans), and membranous nephropathy (most common in Caucasian adults)
  \item Secondary etiologies include diabetic nephropathy, amyloidosis, and lupus nephritis.
\end{itemize}
\section{Risk Factors}


\begin{itemize}[noitemsep]
  \item Genetic predisposition and family history
  \item Advancing age
  \item Lifestyle factors (diet, exercise, smoking, alcohol)
  \item Environmental exposures
  \item Underlying medical conditions
  \item Medication use
\end{itemize}
\section{Signs and Symptoms}

\subsection{Common symptoms}
\begin{itemize}[noitemsep]
  \item You may experience the hallmark tetrad: heavy proteinuria ($>3.5$ g/24 hours or spot urine protein-to-creatinine ratio $>3.5$), hypoalbuminemia ($<3.0$ g/dL), generalized peripheral and periorbital swelling (anasarca due to decreased oncotic pressure and secondary sodium retention), and hyperlipidemia with lipiduria ('oval fat bodies' and Maltese cross appearance under polarized light).

  \item Pain or discomfort in the affected area
  \end{itemize}

\subsection{Emergency warning signs}
\begin{itemize}[noitemsep]
  \item Severe or worsening symptoms of nephrotic syndrome require immediate medical evaluation
\end{itemize}
\section{Progression and Stages}
The condition may progress through different stages of severity over time. Early stages often involve mild or subtle symptoms that may be overlooked. Moderate stages involve more noticeable symptoms affecting daily function. Advanced stages may involve significant impairment and complications.
\section{Complications}
Complications include deep vein blood clot formation/lung blockage by a blood clot (due to urinary loss of antithrombin III) and infection (loss of immunoglobulins).

\begin{itemize}[noitemsep]
  \item Disease progression leading to worsening symptoms
  \item Secondary infections or injuries
  \item Side effects of treatment
  \item Reduced quality of life
  \item Emotional and psychological impact
\end{itemize}
\section{Diagnosis}
Doctors diagnose this using 24-hour urine collection or spot ratio, serum chemistries, and kidney biopsy.

\begin{itemize}[noitemsep]
  \item Comprehensive medical history and physical examination
  \item Laboratory tests to assess organ function
  \item Imaging studies to visualize affected structures
  \item Specialized testing depending on the affected system
  \item Consultation with relevant specialists
\end{itemize}
\section{Prevention}
\begin{itemize}[noitemsep]
  \item Maintain a healthy lifestyle with balanced diet and exercise
  \item Avoid known risk factors
  \item Attend regular medical check-ups for early detection
  \item Follow recommended screening guidelines
\end{itemize}
\section{Treatment Options}
Treatment options include ACE inhibitors/ARBs (to reduce intraglomerular pressure), loop diuretics, statins, anticoagulation when indicated, and systemic corticosteroids or immunosuppressive agents. Symptomatic management and supportive care Enzyme replacement therapy for certain conditions Gene therapy (emerging for select conditions) Dietary modifications (e.g., phenylketonuria) Regular monitoring for associated complications Physical and occupational therapy Genetic counseling for family planning Participation in clinical trials for new therapies

\begin{itemize}[noitemsep]
  \item Medications to manage symptoms and address underlying mechanisms
  \item Lifestyle modifications and self-care measures
  \item Physical therapy and rehabilitation
  \item Surgical intervention when indicated
  \item Regular monitoring and follow-up care
\end{itemize}
\section{Living with It}
\begin{itemize}[noitemsep]
  \item Follow your treatment plan as prescribed
  \item Maintain regular follow-up with your healthcare provider
  \item Make necessary lifestyle adjustments
  \item Seek support from family, friends, or support groups
  \item Stay informed about your condition
\end{itemize}
\section{Outlook}
Recovery varies by underlying cause.
